\begin{thebibliography}{99}
\footnotesize\itemsep=0pt

\bibitem{AnderssonDahl1}
Andersson L., Dahl M., Scalar curvature rigidity for asymptotically locally
 hyperbolic manifolds, \href{https://doi.org/10.1023/A:1006547905892}{\textit{Ann. Global Anal. Geom.}} \textbf{16} (1998),
 1--27, \href{https://arxiv.org/abs/dg-ga/9707017}{arXiv:dg-ga/9707017}.

\bibitem{Bar3}
B\"ar C., Lower eigenvalue estimates for {D}irac operators, \href{https://doi.org/10.1007/BF01444701}{\textit{Math. Ann.}}
 \textbf{293} (1992), 39--46.

\bibitem{Bar1}
B\"ar C., Extrinsic bounds for eigenvalues of the {D}irac operator,
 \href{https://doi.org/10.1023/A:1006550532236}{\textit{Ann. Global Anal. Geom.}} \textbf{16} (1998), 573--596.

\bibitem{Bartnik1}
Bartnik R., The mass of an asymptotically flat manifold, \href{https://doi.org/10.1002/cpa.3160390505}{\textit{Comm. Pure
 Appl. Math.}} \textbf{39} (1986), 661--693.

\bibitem{Bartnik4}
Bartnik R., Quasi-spherical metrics and prescribed scalar curvature,
 \href{https://doi.org/10.4310/jdg/1214453422}{\textit{J.~Differential Geom.}} \textbf{37} (1993), 31--71.

\bibitem{BartnikChrusciel}
Bartnik R.A., Chru\'sciel P.T., Boundary value problems for {D}irac-type
 equations, \href{https://doi.org/10.1515/crll.2005.2005.579.13}{\textit{J.~Reine Angew. Math.}} \textbf{579} (2005), 13--73,
 \href{https://arxiv.org/abs/math.DG/0307278}{arXiv:math.DG/0307278}.

\bibitem{BourguignonHijaziMilhoratMoroianu}
Bourguignon J.-P., Hijazi O., Milhorat J.L., Moroianu A., Moroianu S., A
 spinorial approach to {R}iemannian and conformal geometry, \textit{EMS Monogr. Math.},
 \href{https://doi.org/10.4171/136}{Eur. Math. Soc. (EMS)}, Z\"urich, 2015.

\bibitem{Bray}
Bray H.L., Proof of the {R}iemannian {P}enrose inequality using the positive
 mass theorem, \href{https://doi.org/10.4310/jdg/1090349428}{\textit{J.~Differential Geom.}} \textbf{59} (2001), 177--267,
 \href{https://arxiv.org/abs/math.DG/9911173}{arXiv:math.DG/9911173}.

\bibitem{BrendleHungWang}
Brendle S., Hung P.-K., Wang M.-T., A {M}inkowski inequality for hypersurfaces in
 the anti-de {S}itter--{S}chwarzschild manifold, \href{https://doi.org/10.1002/cpa.21556}{\textit{Comm. Pure Appl.
 Math.}} \textbf{69} (2016), 124--144, \href{https://arxiv.org/abs/1209.0669}{arXiv:1209.0669}.

\bibitem{Chrusciel1}
Chru\'sciel P., Boundary conditions at spatial infinity from a {H}amiltonian
 point of view, in Topological Properties and Global Structure of Space-Time
 ({E}rice, 1985), \textit{NATO Adv. Sci. Inst. Ser.~B Phys.}, Vol.~138,
 \href{https://doi.org/10.1007/978-1-4899-3626-4_5}{Plenum}, New York, 1986, 49--59.

\bibitem{ChruscielHerzlich}
Chru\'sciel P.T., Herzlich M., The mass of asymptotically hyperbolic
 {R}iemannian manifolds, \href{https://doi.org/10.2140/pjm.2003.212.231}{\textit{Pacific~J. Math.}} \textbf{212} (2003),
 231--264, \href{https://arxiv.org/abs/math.DG/0110035}{arXiv:math.DG/0110035}.

\bibitem{EichmairMiaoWang}
Eichmair M., Miao P., Wang X., Extension of a theorem of {S}hi and {T}am,
 \href{https://doi.org/10.1007/s00526-011-0402-2}{\textit{Calc. Var. Partial Differential Equations}} \textbf{43} (2012),
 45--56, \href{https://arxiv.org/abs/0911.0377}{arXiv:0911.0377}.

\bibitem{FanShiTam}
Fan X.-Q., Shi Y., Tam L.-F., Large-sphere and small-sphere limits of the
 {B}rown--{Y}ork mass, \href{https://doi.org/10.4310/CAG.2009.v17.n1.a3}{\textit{Comm. Anal. Geom.}} \textbf{17} (2009), 37--72,
 \href{https://arxiv.org/abs/0711.2552}{arXiv:0711.2552}.

\bibitem{Friedrich3}
Friedrich Th., Der erste {E}igenwert des {D}irac-{O}perators einer kompakten,
 {R}iemannschen {M}annigfaltigkeit nichtnegativer {S}kalarkr\"ummung,
 \href{https://doi.org/10.1002/mana.19800970111}{\textit{Math. Nachr.}} \textbf{97} (1980), 117--146.

\bibitem{GeWangWu}
Ge Y., Wang G., Wu J., Hyperbolic {A}lexandrov--{F}enchel quermassintegral
 inequalities~{II}, \href{https://doi.org/10.4310/jdg/1406552250}{\textit{J.~Differential Geom.}} \textbf{98} (2014),
 237--260, \href{https://arxiv.org/abs/1304.1417}{arXiv:1304.1417}.

\bibitem{Ginoux2}
Ginoux N., Une nouvelle estimation extrins\`eque du spectre de l'op\'erateur de
 {D}irac, \href{https://doi.org/10.1016/S1631-073X(03)00206-1}{\textit{C.~R.~Math. Acad. Sci. Paris}} \textbf{336} (2003), 829--832.

\bibitem{Herzlich1}
Herzlich M., A {P}enrose-like inequality for the mass of {R}iemannian
 asymptotically flat manifolds, \href{https://doi.org/10.1007/s002200050159}{\textit{Comm. Math. Phys.}} \textbf{188}
 (1997), 121--133.

\bibitem{Herzlich2}
Herzlich M., Minimal surfaces, the {D}irac operator and the {P}enrose
 inequality, in S\'eminaire de {T}h\'eorie {S}pectrale et {G}\'eom\'etrie, {A}nn\'ee 2001--2002, \textit{S\'emin. Th\'eor. Spectr. G\'eom.}, Vol.~20,
 Universit\'e de Grenoble I, Saint-Martin-d'H\`eres, 2002, 9--16.

\bibitem{Hijazi2}
Hijazi O., A conformal lower bound for the smallest eigenvalue of the {D}irac
 operator and {K}illing spinors, \href{https://doi.org/10.1007/BF01210797}{\textit{Comm. Math. Phys.}} \textbf{104}
 (1986), 151--162.

\bibitem{Hijazi1}
Hijazi O., Premi\`ere valeur propre de l'op\'erateur de {D}irac et nombre de
 {Y}amabe, \textit{C.~R.~Acad. Sci. Paris S\'er.~I Math.} \textbf{313} (1991),
 865--868.

\bibitem{HijaziMontielRaulot3}
Hijazi O., Montiel S., Raulot S., A holographic principle for the existence of
 imaginary {K}illing spinors, \href{https://doi.org/10.1016/j.geomphys.2015.01.012}{\textit{J.~Geom. Phys.}} \textbf{91} (2015),
 12--28, \href{https://arxiv.org/abs/1502.04091}{arXiv:1502.04091}.

\bibitem{HijaziMontielRaulot4}
Hijazi O., Montiel S., Raulot S., A positive mass theorem for asymptotically
 hyperbolic manifolds with inner boundary, \href{https://doi.org/10.1142/S0129167X15501013}{\textit{Internat.~J.~Math.}}
 \textbf{26} (2015), 1550101, 17~pages.

\bibitem{HijaziMontielRoldan}
Hijazi O., Montiel S., Rold\'an A., Dirac operators on hypersurfaces of
 manifolds with negative scalar curvature, \href{https://doi.org/10.1023/A:1022808916165}{\textit{Ann. Global Anal. Geom.}}
 \textbf{23} (2003), 247--264.

\bibitem{HijaziMontielZhang1}
Hijazi O., Montiel S., Zhang X., Dirac operator on embedded hypersurfaces,
 \href{https://doi.org/10.4310/MRL.2001.v8.n2.a8}{\textit{Math. Res. Lett.}} \textbf{8} (2001), 195--208,
 \href{https://arxiv.org/abs/math.DG/0012262}{arXiv:math.DG/0012262}.

\bibitem{Kwong1}
Kwong K.-K., On the positivity of a quasi-local mass in general dimensions,
 \href{https://doi.org/10.4310/CAG.2013.v21.n4.a3}{\textit{Comm. Anal. Geom.}} \textbf{21} (2013), 847--871, \href{https://arxiv.org/abs/1207.7333}{arXiv:1207.7333}.

\bibitem{Nirenberg}
Nirenberg L., The {W}eyl and {M}inkowski problems in differential geometry in
 the large, \href{https://doi.org/10.1002/cpa.3160060303}{\textit{Comm. Pure Appl. Math.}} \textbf{6} (1953), 337--394.

\bibitem{Pogorelov}
Pogorelov A.V., Regularity of a convex surface with given {G}aussian curvature,
 \textit{Mat. Sb.} \textbf{31} (1952), 88--103.

\bibitem{SchoenYau5}
Schoen R., Yau S.-T., Proof of the positive mass theorem.~{II}, \href{https://doi.org/10.1007/BF01942062}{\textit{Comm.
 Math. Phys.}} \textbf{79} (1981), 231--260.

\bibitem{SchoenYau3}
Schoen R., Yau S.-T., Positive scalar curvature and minimal hypersurface
 singularities, in Surveys in Differential Geometry 2019, {D}ifferential
 geometry, {C}alabi--{Y}au Theory, and General Relativity. {P}art~2,
 \textit{Surv. Differ. Geom.}, Vol.~24, \href{https://dx.doi.org/10.4310/SDG.2019.v24.n1.a10}{Int. Press}, Boston, MA, 2019,
 441--480, \href{https://arxiv.org/abs/1704.05490}{arXiv:1704.05490}.

\bibitem{ShiTam1}
Shi Y., Tam L.-F., Positive mass theorem and the boundary behaviors of compact
 manifolds with nonnegative scalar curvature, \href{https://doi.org/10.4310/jdg/1090425530}{\textit{J.~Differential Geom.}}
 \textbf{62} (2002), 79--125, \href{https://arxiv.org/abs/math.DG/0301047}{arXiv:math.DG/0301047}.

\bibitem{ShiTam2}
Shi Y., Tam L.-F., Rigidity of compact manifolds and positivity of quasi-local
 mass, \href{https://doi.org/10.1088/0264-9381/24/9/013}{\textit{Classical Quantum Gravity}} \textbf{24} (2007), 2357--2366.

\bibitem{WangYau0}
Wang M.T., Yau S.-T., A generalization of {L}iu--{Y}au's quasi-local mass,
 \href{https://dx.doi.org/10.4310/CAG.2007.v15.n2.a2}{\textit{Comm. Anal. Geom.}} \textbf{15} (2007), 249--282,
 \href{https://arxiv.org/abs/math.DG/0602321}{arXiv:math.DG/0602321}.

\bibitem{Witten1}
Witten E., A new proof of the positive energy theorem, \href{https://doi.org/10.1007/BF01208277}{\textit{Comm. Math.
 Phys.}} \textbf{80} (1981), 381--402.

\end{thebibliography}
